\section{Configuration and schema reference}
\label{app:config}

\subsection{Minimal economic-geography configuration}

\begin{lstlisting}
schema_version = 1
name = "my-network"

[input]
adapter = "generic_csv_v1"
nodes = "data/nodes.csv"
edge_modes = "data/edge_modes.csv"

[model]
spatial_specification = "economic_geography"
alpha = 0.10
beta = -0.30
sigma = 9.0
modal_specification = "choice_logsum"
eta = 1.099
route_curvature = "theorem"

[congestion]
specification = "edge"
[congestion.edge]
road = 0.092

[policy]
mode = "road"
unit = "both"
shock_fraction = 0.01

[output]
directory = "output"

[diagnostics]
tolerance = 1.0e-10
condition_limit = 1.0e12
\end{lstlisting}

\subsection{Spatial specification}

\begin{description}
  \item[\code{economic\_geography}.] Requires \(\alpha,\beta,\sigma\). Route
  curvature is \(\sigma-1\).
  \item[\code{urban\_commuting}.] Requires \(\alpha,\beta,\theta\). Internally
  \(\sigma=\theta+1\) supplies the shared route curvature. Residence and
  workplace columns replace labor and income shares.
\end{description}

\subsection{Modal specification}

\begin{description}
  \item[\code{choice\_logsum}.] Negative-power modal cost index. This is the
  active paper and urban extension.
  \item[\code{component\_ces}.] Positive-power legacy component index. Use only
  for declared provenance replications.
\end{description}
Both require finite \(\eta>0\).

\subsection{Congestion specification}

\subsubsection*{None}
\begin{lstlisting}
[congestion]
specification = "none"
\end{lstlisting}

\subsubsection*{Mode-level edge congestion}
\begin{lstlisting}
[congestion]
specification = "edge"
[congestion.edge]
road = 0.092
\end{lstlisting}

\subsubsection*{Edge-specific congestion}
\begin{lstlisting}
[congestion]
specification = "edge"
source = "input_column"
column = "congestion_elasticity"
scale = 1.0
\end{lstlisting}
Every active edge-mode row must contain a finite nonnegative value in the
selected column. Mode-level and column-based values cannot be mixed.

\subsubsection*{Endpoint-terminal congestion}
\begin{lstlisting}
[congestion]
specification = "endpoint_terminal"
endpoint_scale = 1.0
[congestion.terminal]
rail = 0.096
\end{lstlisting}
Every positive-congestion pair must identify both endpoint terminals.

\subsubsection*{Composite congestion}
\begin{lstlisting}
[congestion]
specification = "composite"
endpoint_scale = 1.0
[congestion.edge]
road = 0.092
[congestion.terminal]
transit = 0.04
\end{lstlisting}

\subsection{Policy units}

\begin{description}
  \item[\code{directed\_arc}.] Write directed results only.
  \item[\code{physical\_link}.] Require reciprocal pairs and write physical
  results only.
  \item[\code{both}.] Write both result sets.
\end{description}

\subsection{Economic-geography node schema}

\begin{longtable}{p{0.24\textwidth}p{0.16\textwidth}p{0.50\textwidth}}
\toprule
Column & Required & Meaning \\
\midrule
\endhead
\code{node\_id} & Yes & Stable unique identifier \\
\code{labor} & Yes & Positive normalized labor share \\
\code{income} & Yes & Positive normalized income share \\
\code{longitude} & No & Finite map coordinate \\
\code{latitude} & No & Finite map coordinate \\
\bottomrule
\end{longtable}

\subsection{Urban node schema}

\begin{longtable}{p{0.24\textwidth}p{0.16\textwidth}p{0.50\textwidth}}
\toprule
Column & Required & Meaning \\
\midrule
\endhead
\code{node\_id} & Yes & Stable unique identifier \\
\code{residents} & Yes & Normalized residence margin \(l^R\) \\
\code{employment} & Yes & Normalized workplace margin \(l^F\) \\
\code{longitude} & No & Finite map coordinate \\
\code{latitude} & No & Finite map coordinate \\
\bottomrule
\end{longtable}

\subsection{Edge-mode schema}

\begin{longtable}{p{0.27\textwidth}p{0.14\textwidth}p{0.48\textwidth}}
\toprule
Column & Required & Meaning \\
\midrule
\endhead
\code{edge\_id} & Yes & Directed edge identifier, common across modal rows
when appropriate \\
\code{physical\_link\_id} & Yes & Identifier used for reciprocal aggregation \\
\code{origin} & Yes & Existing node ID \\
\code{destination} & Yes & Existing node ID, distinct from origin \\
\code{mode} & Yes & Nonempty mode name \\
\code{flow} & Yes & Finite positive normalized value flow \\
\code{origin\_terminal\_id} & Conditional & Required for endpoint-terminal
congestion \\
\code{destination\_}\newline\code{terminal\_id} & Conditional & Required for
endpoint-terminal congestion \\
\code{congestion\_elasticity} & Conditional & Required when selected as the
edge-congestion source \\
\bottomrule
\end{longtable}

\subsection{Sensitivity}

\begin{lstlisting}
[sensitivity]
alpha = [0.06, 0.08, 0.10, 0.12, 0.14]
beta = [-0.38, -0.34, -0.30, -0.26, -0.22]
eta = [0.75, 0.90, 1.099, 1.25, 1.40]
lambda_road = [0.00, 0.04, 0.092, 0.14]
\end{lstlisting}
Column-based congestion uses \code{edge\_congestion\_scale} rather than a
mode-specific lambda path.
